% Tracing, on the AST of a small function: R evaluates the grey nodes at
% trace time and they leave no trace; the blue nodes reach an overloaded
% array operation and are transcribed into the graph as primitives.
% intercept-ast.tex sets \intercepting to 1 and draws the same tree as an
% interpreter-intercepting compiler sees it: it has to execute every node.
\providecommand{\intercepting}{0}
\documentclass[tikz, border=8pt]{standalone}
\usepackage{forest}
\usetikzlibrary{arrows.meta, positioning, fit, backgrounds, calc}

\begin{document}
\ifnum\intercepting=1
  \forestset{staged/.style={draw=orange!75!black, fill=orange!10, thick, text=black}}
\else
  \forestset{staged/.style={draw=black!35, dashed, fill=black!4, text=black!55}}
\fi
\begin{forest}
  for tree={
    draw, rounded corners=2pt, thick, font=\ttfamily\small,
    minimum height=0.6cm, minimum width=0.8cm, inner sep=3pt,
    s sep=5mm, l sep=7mm, edge={thick, draw=black!45}
  },
  where level=0{s sep=16mm}{},
  prim/.style={draw=blue!65, fill=blue!12, very thick},
  input/.style={draw=black!65, fill=white},
  [\{, staged, name=root
    [<-, staged
      [s, staged]
      [*, staged, name=pi
        [2, staged, name=two] [pi, staged, name=pisym]]]
    [<-, staged
      [y, staged]
      [\%*\%, prim, name=mm
        [x, input, name=xin] [w, input, name=win]]]
    [exp, prim, name=exp
      [*, prim, name=mul
        [y, staged] [s, staged]]]
  ]
  \tikzset{
    note/.style={font=\footnotesize, align=left, inner sep=2pt},
    gnote/.style={note, text=black!55},
    bnote/.style={note, text=blue!65!black}
  }
  \ifnum\intercepting=1
    \node[note, align=center, text=orange!60!black] at ([yshift=-0.55cm]$(two.south)!0.5!(pisym.south)$) {folded by the compiler:\\[-1pt]it must know \texttt{*}, \texttt{pi}};
    \node[note, align=left, text=orange!60!black, right=0.3cm of root] {anything it cannot\\[-1pt]model: a graph break};
  \else
    \node[gnote, align=center] at ([yshift=-0.55cm]$(two.south)!0.5!(pisym.south)$) {evaluated by R:\\[-1pt]\texttt{6.283185}};
  \fi
  \node[bnote, align=center] at ([yshift=-0.4cm]$(xin.south)!0.5!(win.south)$) {\texttt{\%1 = dot\_general}};
  \node[bnote, right=0.15cm of exp] {\texttt{\%4 = exp}};
  \node[bnote, right=0.15cm of mul, align=left] {\texttt{\%2 = broadcast\_in\_axes}\\[-1pt]\texttt{\%3 = mul}};
  % The function, left of the tree.
  \node[draw=black!50, rounded corners=3pt, thick, fill=white, align=left,
        font=\ttfamily\small, inner sep=6pt, anchor=north east]
    (code) at ([xshift=-2.6cm, yshift=0.3cm]root.north -| current bounding box.west)
    {f <- function(x, w) \{\\
     \ \ s <- 2 * pi\\
     \ \ y <- x \%*\% w\\
     \ \ exp(y * s)\\
     \}};
  % The graph it traces to, below the function.
  \node[draw=blue!65, rounded corners=3pt, very thick, fill=blue!5, align=left,
        font=\ttfamily\small, inner sep=6pt, below=0.9cm of code.south west, anchor=north west]
    (graph)
    {\%1 = dot\_general(\%x1, \%x2)\\
     \%2 = broadcast\_in\_axes(6.283185)\\
     \%3 = mul(\%1, \%2)\\
     \%4 = exp(\%3)\\
     return \%4};
  \node[font=\footnotesize\bfseries, above=1pt of code.north west, anchor=south west] {the function};
  \node[font=\footnotesize\bfseries, above=1pt of graph.north west, anchor=south west, text=blue!65!black] {the traced graph};
  \node[font=\footnotesize\bfseries, above=4pt of root] {its AST};
  % Legend.
  \ifnum\intercepting=1
    \node[draw=orange!75!black, fill=orange!10, thick, rounded corners=2pt, minimum width=0.5cm, minimum height=0.35cm,
          below=0.6cm of graph.south west, anchor=north west] (l1) {};
    \node[note, right=0.15cm of l1] {executed symbolically by the compiler: it has to model R's semantics};
  \else
    \node[draw=black!35, dashed, fill=black!4, rounded corners=2pt, minimum width=0.5cm, minimum height=0.35cm,
          below=0.6cm of graph.south west, anchor=north west] (l1) {};
    \node[note, right=0.15cm of l1] {staged out: R evaluates it at trace time, nothing is recorded};
  \fi
  \node[draw=blue!65, fill=blue!12, very thick, rounded corners=2pt, minimum width=0.5cm, minimum height=0.35cm,
        below=0.25cm of l1] (l2) {};
  \node[note, right=0.15cm of l2] {array operation on a \texttt{GraphBox}: transcribed into the graph as primitives};
\end{forest}
\end{document}
